\documentclass{article}

\usepackage[utf8]{inputenc}
\usepackage[german]{babel}
\usepackage{listings}
\usepackage{xcolor}

\lstset{
    basicstyle=\ttfamily\small,
    breaklines=true,
    backgroundcolor=\color{gray!10}
}

\begin{document}
    \begin{center}
        \Large \textbf{Huntsman Challenge 1: Inventarverwaltung und Hostentdeckung} \\ [6pt]
        \normalsize
        Autor: Maximilian Jakob Maag B.Sc. \\
        Version: 1.0
    \end{center}

    \section{Intro}
    In dieser Challenge geht es um die Verwaltung von Infrastruktur durch ein gut organisiertes Ansible Inventory.
    Sie werden lernen, wie man eine große Anzahl von Servern strukturiert erfasst und verwaltet.
    Ihr Ziel ist es, alle von Ihrem Rabbit-Konkurrenten deployten Server zu dokumentieren und zu organisieren.

    \section{Aufgaben}

    \subsection{Aufgabe 1}
    Ihr Rabbit-Konkurrent hat begonnen, Server zu deployen. Erhalten Sie von ihm eine Liste aller deployten Server mit ihren IP-Adressen und Hostnamen.
    \\ \\
    \textbf{Hinweis:} Die Server folgen dem Namenskonzept: \texttt{rabbitname-servernumber.simple-test.org}

    \subsection{Aufgabe 2}
    Erstellen Sie ein strukturiertes Ansible Inventory in YAML Format, das folgende Informationen enthält:
    \begin{itemize}
        \item Alle Hosts mit Hostnamen und IP-Adressen
        \item Gruppierung nach Region (z.B. US-east1, EU-west1, ASIA-east1)
        \item Gruppierung nach Funktionalität (z.B. webserver, database, cache)
        \item Gruppierung nach Betriebssystem Version
    \end{itemize}

    \textbf{Beispiel Inventarstruktur:}
    \begin{lstlisting}[language=yaml]
all:
  children:
    region_us_east:
      hosts:
        rabbit-web01.simple-test.org:
          ansible_host: 192.0.2.1
          function: webserver
    region_eu_west:
      hosts:
        rabbit-web02.simple-test.org:
          ansible_host: 192.0.2.2
          function: webserver
    \end{lstlisting}

    \subsection{Aufgabe 3}
    Implementieren Sie ein Ansible Playbook, das alle verwalteten Hosts auf Erreichbarkeit überprüft.
    Das Playbook soll sich bei jedem Host verbinden und grundlegende Informationen wie OS, Kernel Version und verfügbare Ressourcen erfassen.

    \subsection{Aufgabe 4}
    Erstellen Sie ein dynamisches Inventory Script (oder verwenden Sie ein vorhandenes Ansible Plugin),
    das die Server automatisch aus der Linode API abruft.
    Dies ermöglicht es Ihnen, automatisch neue vom Rabbit deployten Server zu entdecken.

    \subsection{Aufgabe 5}
    Entwickeln Sie ein Namenskonzept für verwaltete Server, das folgende Aspekte berücksichtigt:
    \begin{itemize}
        \item Region oder Standort
        \item Funktionalität oder Service (z.B. web, db, cache)
        \item Instanznummer
        \item Umgebung (production, staging, development)
    \end{itemize}
    
    Dokumentieren Sie dieses Konzept und erstellen Sie ein Playbook, das Host-Metadaten basierend auf dem Namen automatisch generiert.

    \subsection{Aufgabe 6}
    Erstellen Sie eine Dokumentation Ihres Inventars in Markdown Format.
    Diese Dokumentation soll einem neuen Admin ermöglichen, schnell die Struktur und den Aufbau der Infrastruktur zu verstehen.
    \\ \\
    Die Dokumentation sollte enthalten:
    \begin{itemize}
        \item Übersicht aller Hostgruppen
        \item Erklärung des Namenskonzeptes
        \item Beispiele für häufig verwendete Ansible Commands
        \item Troubleshooting Guide
    \end{itemize}

\end{document}
